\documentclass[12pt,a4paper]{article}

\usepackage[T1]{fontenc}
\usepackage[utf8]{inputenc}
\usepackage[romanian]{babel}
\usepackage{geometry}
\usepackage{amsmath}
\usepackage{listings}
\usepackage{xcolor}

\geometry{margin=2.5cm}

\title{Problema Diamantului în C++}
\author{}
\date{}

\begin{document}

\maketitle

\section{Introducere}

Problema diamantului (Diamond Problem) reprezintă una dintre cele mai cunoscute dificultăți asociate moștenirii multiple.

Aceasta apare atunci când o clasă moștenește indirect aceeași clasă de bază prin mai multe căi.

Denumirea provine de la forma ierarhiei de clase:

\begin{verbatim}
        A
       / \
      B   C
       \ /
        D
\end{verbatim}

\section{Motivație}

Considerăm o clasă de bază A.

Clasele B și C moștenesc clasa A.

Clasa D moștenește atât B cât și C.

Astfel, D primește moștenirea lui A prin două căi diferite.

\section{Exemplu clasic}

\begin{verbatim}
class A {
public:
    int x;
};

class B : public A {
};

class C : public A {
};

class D : public B, public C {
};
\end{verbatim}

Grafic:

\begin{verbatim}
        A
       / \
      B   C
       \ /
        D
\end{verbatim}

\section{Problema}

În clasa D există două copii ale clasei A:

\begin{itemize}
\item una moștenită prin B;
\item una moștenită prin C.
\end{itemize}

Astfel:

\begin{verbatim}
D d;

d.x = 10;
\end{verbatim}

produce eroare de compilare.

Compilatorul nu știe la care membru x ne referim.

\section{Ambiguitatea}

Compilatorul observă două posibilități:

\begin{verbatim}
d.B::x
\end{verbatim}

sau

\begin{verbatim}
d.C::x
\end{verbatim}

Deoarece existența celor două căi produce ambiguitate, accesul direct nu este permis.

\section{Demonstrație}

\begin{verbatim}
#include <iostream>

class A {
public:
    int x = 10;
};

class B : public A {
};

class C : public A {
};

class D : public B, public C {
};

int main() {

    D d;

    d.x = 100;

    return 0;
}
\end{verbatim}

Compilatorul va genera eroare deoarece membrul x există de două ori.

\section{Accesarea explicită}

Ambiguitatea poate fi rezolvată temporar folosind operatorul de rezoluție.

\begin{verbatim}
D d;

d.B::x = 10;

d.C::x = 20;
\end{verbatim}

În acest caz specificăm explicit copia utilizată.

\section{De ce este o problemă?}

Problema nu apare doar pentru atribute.

Ea apare și pentru:

\begin{itemize}
\item metode;
\item constructori;
\item variabile membre;
\item orice alt membru moștenit.
\end{itemize}

Exemplu:

\begin{verbatim}
class A {
public:
    void f() {
    }
};
\end{verbatim}

În D vor exista două versiuni ale metodei f().

\section{Soluția - Moștenirea Virtuală}

C++ oferă o soluție numită moștenire virtuală.

Ideea este ca toate clasele derivate să împartă aceeași copie a clasei de bază.

\section{Sintaxa moștenirii virtuale}

\begin{verbatim}
class A {
public:
    int x;
};

class B : virtual public A {
};

class C : virtual public A {
};

class D : public B, public C {
};
\end{verbatim}

Observăm utilizarea cuvântului cheie:

\begin{verbatim}
virtual
\end{verbatim}

în declarațiile claselor B și C.

\section{Rezultatul}

Acum există o singură copie a clasei A.

Grafic:

\begin{verbatim}
          A
         / \
        /   \
       B     C
        \   /
         \ /
          D
\end{verbatim}

Din punct de vedere logic avem aceeași structură, dar în memorie există o singură instanță a lui A.

\section{Exemplu complet}

\begin{verbatim}
#include <iostream>

class A {
public:
    int x;
};

class B : virtual public A {
};

class C : virtual public A {
};

class D : public B, public C {
};

int main() {

    D d;

    d.x = 42;

    std::cout << d.x;

    return 0;
}
\end{verbatim}

Programul compilează și funcționează corect.

\section{Avantaje}

Moștenirea virtuală:

\begin{itemize}
\item elimină ambiguitățile;
\item elimină copiile multiple ale bazei;
\item simplifică ierarhiile complexe;
\item este soluția standard pentru Diamond Problem.
\end{itemize}

\section{Dezavantaje}

\begin{itemize}
\item implementare mai complexă;
\item cost suplimentar la compilare;
\item ierarhiile devin mai greu de urmărit.
\end{itemize}

\section{Recapitulare}

Problema diamantului apare atunci când aceeași clasă de bază este moștenită prin mai multe căi.

\begin{verbatim}
        A
       / \
      B   C
       \ /
        D
\end{verbatim}

Consecința este existența a două copii ale clasei A în D.

Soluția este utilizarea moștenirii virtuale:

\begin{verbatim}
class B : virtual public A { };

class C : virtual public A { };
\end{verbatim}

Astfel, toate clasele împart aceeași copie a clasei de bază și ambiguitatea dispare.

\end{document}